% Part: first-order-logic
% Chapter: sequent-calculus
% Section: proof-theoretic-notions

\documentclass[../../../include/open-logic-section]{subfiles}

\begin{document}

\begin{editorial}
  دا برخه د سېکوېنټ حساب د اشتقاق وړتيا د اړيکې او سازګارۍ
  تعريفونه راغونډوي۔
  \textbf{د ژباړې د نسخې سمون (OLSQ-003):} په سرچينه کښې د طبيعي
  استنتاج نوم راغلے و؛ دلته د څپرکي او ورپسې تعريفونو له مخې د
  سېکوېنټ حساب نوم کارول کېږي۔
\end{editorial}

\iftag{FOL}
      {\olfileid{fol}{seq}{ptn}}
      {\olfileid{pl}{seq}{ptn}}

\olsection{د ثبوت-تيوريکي مفاهيم}

\begin{explain}
لکه څنګه چې مو يو شمېر مهم معنايي مفاهيم (اعتبار، معنايي استلزام
او د صدق وړتيا) تعريف کړي دي، اوس ورسره سم \emph{ثبوت-تيوريکي
مفاهيم} تعريفوو۔ دا مفاهيم په !!{structure}s کښې د !!{sentence}s
د صدق پر بنسټ نۀ، بلکې د ځينو سېکوېنټونو د !!{derivability} يا
!!{nonderivability} پر بنسټ تعريفېږي۔ دا يوه مهمه موندنه وه چې دغه
مفاهيم سره سمون خوري۔ د دې سمون منځپانګه د \emph{سموالي} او
\emph{بشپړتيا قضيه} ده۔
\end{explain}

\begin{defn}[قضيې]
!!^a{sentence}~$!A$ هله \emph{قضيه} ده چې په~$\Log{LK}$ کښې د
سېکوېنټ $\quad \Sequent !A$ !!a{derivation} موجود وي۔ که $!A$
قضيه وي، $\Proves !A$ ليکو؛ او که نۀ وي، $\Proves/ !A$ ليکو۔
\end{defn}

\begin{defn}[!!^{derivability}]
!!^a{sentence}~$!A$ د !!{sentence}s له سټ~$\Gamma$ څخه
\emph{!!{derivable}} ده، يعنې $\Gamma \Proves !A$، هله او يوازې
هله چې يو متناهي فرعي سټ~$\Gamma_0 \subseteq \Gamma$ او په
$\Gamma_0$ کښې د !!{sentence}s يو لړ~$\Gamma_0'$ موجود وي، داسې
چې $\Log{LK}$ د $\Gamma_0' \Sequent !A$ سېکوېنټ !!{derive}s کړي۔
که $!A$ له $\Gamma$ څخه !!{derivable} نۀ وي، $\Gamma \Proves/ !A$
ليکو۔
\end{defn}

د انقباض، کمزوري کولو او تبادلې د قاعدو له امله، په~$\Gamma_0'$
کښې د !!{sentence}s ترتيب او شمېر اهميت نۀ لري: که سېکوېنټ
$\Gamma_0' \Sequent !A$ !!{derivable} وي، نو د هر داسې
$\Gamma_0''$ دپاره $\Gamma_0'' \Sequent !A$ هم د اشتقاق وړ دے
چې د~$\Gamma_0'$ هماغه !!{sentence}s لري۔ مثلاً که
$\Gamma_0 = \{!B, !C\}$ وي، نو هم
$\Gamma_0' = \tuple{!B, !B, !C}$ او هم
$\Gamma_0'' = \tuple{!C, !C, !B}$ داسې لړونه دي چې يوازې د
$\Gamma_0$ !!{sentence}s لري۔ که د يوه لړ لرونکے سېکوېنټ
!!{derivable} وي، د بل لړ لرونکے هم دے، مثلاً:
\begin{prooftree}
  \AxiomC{}
  \Deduce$!B, !B, !C \fCenter !A$
  \RightLabel{\LeftR{\Contraction}}
  \UnaryInf$!B, !C \fCenter !A$
  \RightLabel{\LeftR{\Exchange}}
  \UnaryInf$!C, !B \fCenter !A$
  \RightLabel{\LeftR{\Weakening}}
  \UnaryInf$!C, !C, !B \fCenter !A$
\end{prooftree}
له دې وروسته به وايو چې که $\Gamma_0$ د !!{sentence}s يو متناهي
سټ وي، نو $\Gamma_0 \Sequent !A$ هر هغه سېکوېنټ دے چې مقدم ئې
په~$\Gamma_0$ کښې د !!{sentence}s يو لړ وي؛ او د اړتيا په وخت کښې
انقباضونه، تبادلې او کمزوري کول په ناويلې توګه ورګډوو۔

\begin{defn}[سازګاري]
د جملو سټ~$\Gamma$ هله او يوازې هله \emph{ناسازګار} دے چې يو
متناهي فرعي سټ~$\Gamma_0 \subseteq \Gamma$ موجود وي، داسې چې
$\Log{LK}$ د $\Gamma_0 \Sequent \quad$ سېکوېنټ !!{derive}s کړي۔
که $\Gamma$ ناسازګار نۀ وي، يعنې که د هر متناهي
$\Gamma_0 \subseteq \Gamma$ دپاره $\Log{LK}$ د
$\Gamma_0 \Sequent \quad$ سېکوېنټ !!{derive} نۀ کړي، نو هغه
\emph{سازګار} بولو۔
\end{defn}

\begin{prop}[انعکاسيت]
\ollabel{prop:reflexivity}
که $!A \in \Gamma$ وي، نو $\Gamma \Proves !A$۔
\end{prop}

\begin{proof}
لومړنے سېکوېنټ $!A \Sequent !A$ !!{derivable} دے، او
$\{!A\} \subseteq \Gamma$۔
\end{proof}
  
\begin{prop}[يکنواختي]
\ollabel{prop:monotonicity}
که $\Gamma \subseteq \Delta$ او $\Gamma \Proves !A$ وي، نو
$\Delta \Proves !A$۔
\end{prop}

\begin{proof}
فرض کړئ $\Gamma \Proves !A$، يعنې يو متناهي
$\Gamma_0 \subseteq \Gamma$ شته چې
$\Gamma_0 \Sequent !A$ !!{derivable} دے۔ ځکه چې
$\Gamma \subseteq \Delta$، نو $\Gamma_0$ د~$\Delta$ هم يو متناهي
فرعي سټ دے۔ له دې امله د $\Gamma_0 \Sequent !A$ !!{derivation}
هم $\Delta \Proves !A$ ښيي۔
\end{proof}

\begin{prop}[انتقاليت]
\ollabel{prop:transitivity}
که $\Gamma \Proves !A$ او $\{!A\} \cup \Delta \Proves !B$ وي، نو
$\Gamma \cup \Delta \Proves !B$۔
\end{prop}

\begin{proof}
که $\Gamma \Proves !A$ وي، يو متناهي
$\Gamma_0 \subseteq \Gamma$ او د
$\Gamma_0 \Sequent !A$ !!a{derivation}~$\pi_0$ شته۔ که
$\{!A\} \cup \Delta \Proves !B$ وي، نو د کوم متناهي فرعي سټ
$\Delta_0 \subseteq \Delta$ دپاره د
$!A, \Delta_0 \Sequent !B$ !!a{derivation}~$\pi_1$ شته۔ لاندې
!!{derivation} وګورئ:
\begin{prooftree}
  \AxiomC{}
  \RightLabel{$\pi_0$}
  \Deduce$\Gamma_0 \fCenter !A$
  \AxiomC{}
  \RightLabel{$\pi_1$}
  \Deduce$!A, \Delta_0 \fCenter !B$
  \RightLabel{\Cut}
  \BinaryInf$\Gamma_0, \Delta_0 \fCenter !B$
\end{prooftree}
ځکه چې $\Gamma_0 \cup \Delta_0 \subseteq \Gamma \cup \Delta$، نو
دا $\Gamma \cup \Delta \Proves !B$ ښيي۔
\end{proof}

پام وکړئ، په ځانګړې توګه د دې معنٰی دا ده چې که
$\Gamma \Proves !A$ او $!A \Proves !B$ وي، نو
$\Gamma \Proves !B$۔ همدارنګه ترې دا هم راځي چې که
$!A_1, \dots, !A_n \Proves !B$ وي او د هر~$i$ دپاره
$\Gamma \Proves !A_i$ وي، نو $\Gamma \Proves !B$۔

\begin{prop}
\ollabel{prop:incons}
$\Gamma$ هله او يوازې هله ناسازګار دے چې د هرې جملې~$!A$ دپاره
$\Gamma \Proves {!A}$ وي۔
\end{prop}

\begin{proof}
تمرين۔
\end{proof}

\begin{prob}
\olref[fol][seq][ptn]{prop:incons} ثابت کړئ۔
\end{prob}

\begin{prop}[متناهي شاهد]
  \ollabel{prop:proves-compact}
  \begin{enumerate}
  \item که $\Gamma \Proves !A$ وي، نو يو متناهي فرعي سټ
    $\Gamma_0 \subseteq \Gamma$ شته چې $\Gamma_0 \Proves !A$۔
  \item که د~$\Gamma$ هر متناهي فرعي سټ سازګار وي، نو $\Gamma$
    سازګار دے۔
  \end{enumerate}
\end{prop}

\begin{proof}
  \begin{enumerate}
    \item که $\Gamma \Proves !A$ وي، نو يو متناهي فرعي سټ
      $\Gamma_0 \subseteq \Gamma$ شته، داسې چې د
      $\Gamma_0 \Sequent !A$ سېکوېنټ !!a{derivation} لري۔ له دې امله
      $\Gamma_0 \Proves !A$۔
    \item که $\Gamma$ ناسازګار وي، نو يو متناهي فرعي سټ
      $\Gamma_0 \subseteq \Gamma$ شته، داسې چې $\Log{LK}$ د
      $\Gamma_0 \Sequent \quad$ سېکوېنټ !!{derive}s کړي۔ خو بيا
      $\Gamma_0$ د~$\Gamma$ يو ناسازګار متناهي فرعي سټ دے۔
  \end{enumerate}
\end{proof}

\end{document}
